\documentclass[11pt, ukrainian]{article}
\usepackage[a4paper]{geometry}
\setlength\parindent{0mm}
\geometry{top=1.1cm,bottom=1.1cm,left=1.0cm,right=1.0cm}
\pagestyle{empty}
\usepackage[T2A]{fontenc}
\usepackage[utf8]{inputenc}
\usepackage[ukrainian]{babel}
\usepackage{graphicx}
\begin{document}
%begin
\begin{large}
\textbf{Контрольна робота, варіант \No 1}\par

%beginitem
1. Що буде виведено за виконання фрагменту коду?
%code
if (2 > 2)
    cout << "c";
else
    cout << "v";
%code


%enditem

%beginitem
2. Що буде виведено за виконання фрагменту коду?

%code
#include <iostream>
using namespace std;

bool f(int n){
    cout<<"f";
    return n;
}

int main(){
    cout<<(f(-6) && f(0));
    return 0;
}
%code

%enditem

%beginitem
3. Що буде виведено за виконання фрагменту коду?
%code
#include <iostream>
using namespace std;

int f(int &x, int &y){
    x = 4;
    y+= 3;
    return x;
}

int main(){
    int a = 5, b = 7;
    a = f(a, a);
    cout << a << ":" << b <<':';
    {
        int a = 2, b = 6;
        cout << ((a>6) && ((b+=1) > 6))
             << a << ":" << b << ":";
    }
    {
        inta = 9;
        cout << a << ':';
    }
    cout << a << ":" << b << endl;
    return 0;
}
%code


%enditem
\end{large}
%end
\vspace{20mm}
%begin
\begin{large}
\textbf{Контрольна робота, варіант \No 2}\par

%beginitem
1. Що буде виведено за виконання фрагменту коду?
%code
if (9 >= 9)
    cout << "u";
else
    cout << "o";
%code


%enditem

%beginitem
2. Що буде виведено за виконання фрагменту коду?

%code
#include <iostream>
using namespace std;

bool f(int n){
    cout<<"f";
    return n>-1;
}

int main(){
    cout<<(f(1) || f(6));
    return 0;
}
%code

%enditem

%beginitem
3. Що буде виведено за виконання фрагменту коду?
%code
#include <iostream>
using namespace std;

int f(int &x, int &y){
    x = 1;
    y-= 8;
    return y;
}

int main(){
    int a = 5, b = 1;
    b = f(b, b);
    cout << a << ":" << b <<':';
    {
        int a = 3, b = 9;
        cout << ((b>=5) || ((a-=2) >= 5))
             << a << ":" << b << ":";
    }
    {
        inta = 2;
        cout << a << ':';
    }
    cout << a << ":" << b << endl;
    return 0;
}
%code


%enditem
\end{large}
%end
\newpage
%begin
\begin{large}
\textbf{Контрольна робота, варіант \No 3}\par

%beginitem
1. Що буде виведено за виконання фрагменту коду?
%code
if (0 < 0)
    cout << "e";
else
    cout << "s";
%code


%enditem

%beginitem
2. Що буде виведено за виконання фрагменту коду?

%code
#include <iostream>
using namespace std;

bool f(int n){
    cout<<"f";
    return n>=-2;
}

int main(){
    cout<<(f(-8) and f(-4));
    return 0;
}
%code

%enditem

%beginitem
3. Що буде виведено за виконання фрагменту коду?
%code
#include <iostream>
using namespace std;

int f(int &x, int &y){
    x = 6;
    y-= 8;
    return y;
}

int main(){
    int a = 4, b = 7;
    b = f(b, b);
    cout << a << ":" << b <<':';
    {
        int a = 6, b = 2;
        cout << ((b<7) || ((a-=2) < 7))
             << a << ":" << b << ":";
    }
    {
        inta = 5;
        cout << a << ':';
    }
    cout << a << ":" << b << endl;
    return 0;
}
%code


%enditem
\end{large}
%end
\vspace{20mm}
%begin
\begin{large}
\textbf{Контрольна робота, варіант \No 4}\par

%beginitem
1. Що буде виведено за виконання фрагменту коду?
%code
if (14 != 14)
    cout << "y";
else
    cout << "l";
%code


%enditem

%beginitem
2. Що буде виведено за виконання фрагменту коду?

%code
#include <iostream>
using namespace std;

bool f(int n){
    cout<<"f";
    return n;
}

int main(){
    cout<<(f(-1) or f(9));
    return 0;
}
%code

%enditem

%beginitem
3. Що буде виведено за виконання фрагменту коду?
%code
#include <iostream>
using namespace std;

int f(int &x, int &y){
    x = 7;
    y+= 4;
    return x;
}

int main(){
    int a = 3, b = 9;
    a = f(a, a);
    cout << a << ":" << b <<':';
    {
        int a = 8, b = 1;
        cout << ((b<=3) && ((a+=1) <= 3))
             << a << ":" << b << ":";
    }
    {
        inta = 7;
        cout << a << ':';
    }
    cout << a << ":" << b << endl;
    return 0;
}
%code


%enditem
\end{large}
%end
\newpage
\end{document}

